\documentclass[11pt]{article}
\usepackage{fullpage}
\usepackage{amsfonts}
\usepackage{amssymb}
\usepackage{amsmath}
\usepackage{titling}
\usepackage{color}
\usepackage{graphicx}
\usepackage[top=1.0in, bottom=1.0in, left=1.0in, right=1.0in]{geometry}
\usepackage[font=small, format=hang, up,textfont=normal,up,justification=justified,singlelinecheck=false]{caption}
\captionsetup{width=6cm}
\usepackage{wrapfig}
\usepackage{hyperref}
\usepackage{multicol}
\usepackage{vwcol}
\usepackage{enumitem}
\usepackage{afterpage}

\newcommand\blankpage{%
	\null
	\thispagestyle{empty}%
	\addtocounter{page}{-1}%
	\newpage}


\newcommand{\specialcell}[2][l]{%
  \begin{tabular}[#1]{@{}l@{}}#2\end{tabular}}

\newcommand{\specialcellright}[2][r]{%
  \begin{tabular}[#1]{@{}r@{}}#2\end{tabular}}

\setlength\parindent{12pt}

\newcommand{\years}[1]{\marginnote{\scriptsize #1}}
\newcommand{\bigfont}{\fontsize{18}{18}\selectfont}
\newcommand{\smallfont}{\fontsize{7}{7.2}\selectfont}

\newcommand{\snn}{\sqrt{s_{_{NN}}}}



\begin{document}
%\title{Joseph D. Osborn}
%\maketitle
%\date{}

\begin{center}
\Large\textbf{Joseph (Joe) D. Osborn} \\
\normalsize
osbornjd91@gmail.com \\
(859) 433-8738 \\
https://jdosbo.github.io \\
github.com/osbornjd
\end{center}



%%%%%%%%%%%%%%%%%%%%%%%%%%%%%%%%%%%%%%%%%%%%%%%%%%%%%%%
%%%%%%%															%%%%%%%%%%
%%%%%%					RESEARCH INTERESTS						%%%%%%%%%%
%%%%%%%															%%%%%%%%%%
%%%%%%%%%%%%%%%%%%%%%%%%%%%%%%%%%%%%%%%%%%%%%%%%%%%%%%%
\begin{flushleft}
\LARGE\textbf{Research Interests} \\
\normalsize High-energy experimental nuclear physics, software and computing in high-energy physics.  \\ 
\end{flushleft}

%\begin{flushleft}
%\Large\textbf{Research Overview} \\
%\normalsize Experimental particle physicist with five years experience performing precision measurements of proton and nuclei structure. Possess the skill set necessary to develop complex software frameworks in Python and C++.
%\end{flushleft}








%%%%%%%%%%%%%%%%%%%%%%%%%%%%%%%%%%%%%%%%%%%%%%%%%%%%%%%
%%%%%%%															%%%%%%%%%%
%%%%%%					EDUCATION								%%%%%%%%%%
%%%%%%%															%%%%%%%%%%
%%%%%%%%%%%%%%%%%%%%%%%%%%%%%%%%%%%%%%%%%%%%%%%%%%%%%%%


\vspace{3pt}
\begin{flushleft}
\LARGE\textbf{Education} \\
\normalsize 

University of Michigan, Ph.D. in Physics, 2018. \\
University of Michigan, M.S. in Physics, 2015. \\
University of Kentucky, B.S. in Physics, B.S. in Mathematics, 2013.\\


\end{flushleft}

\vspace{3pt}
\begin{flushleft}
	\Large\textbf{Software and Computational Fluency} \\
	\normalsize
	Languages: C/C++, Java, Python, \LaTeX, Bash \\
	Tools: git, Jenkins, Linux, UML, Doxygen, Docker, Singularity, Serverless/FaaS, pybind\\
	HEP Packages: ROOT, ACTS, GEANT, Fun4All, RooUnfold
\end{flushleft}

\vspace{3pt}

\begin{flushleft}
\LARGE\textbf{Ph.D. Thesis} \\
\normalsize

		Nonperturbative factorization breaking and color entanglement effects in dihadron and direct photon-hadron angular correlations in $p$$+$$p$ and $p$$+$A collisions. J. D. Osborn, University of Michigan, May 29, 2018. hep-ex/1806.07763.
	
\end{flushleft}





%%%%%%%%%%%%%%%%%%%%%%%%%%%%%%%%%%%%%%%%%%%%%%%%%%%%%%%
%%%%%%%															%%%%%%%%%%
%%%%%%					RESEARCH POSITIONS						%%%%%%%%%%
%%%%%%%															%%%%%%%%%%
%%%%%%%%%%%%%%%%%%%%%%%%%%%%%%%%%%%%%%%%%%%%%%%%%%%%%%%




\vspace{3pt}
\begin{flushleft}
\LARGE\textbf{Research Positions} \\
%\vspace{0.2cm}
\normalsize
\vspace{-10pt}
\begin{description}[style=nextline]
	\item[\noindent Advanced Appications Engineer, Brookhaven National Laboratory]
	08/2022-present\\
	Physics Department
	
	\item[\noindent Associate Research Scientist, Oak Ridge National Laboratory] 06/2021-08/2022	\\
	Computing and Computational Sciences Directorate
	
	
	\item [\noindent Postdoctoral Research Associate, Oak Ridge National Laboratory] 07/2019-05/2021 \\
		Computing and Computational Sciences Directorate

	\item [Visiting Scholar, University of Michigan ] 06/2019-06/2020 \\
			Department of Physics 
			
	\item [Postdoctoral Research Fellow, University of Michigan ] 06/2018-06/2019 \\
	Department of Physics

	\item [Graduate Research Assistant, University of Michigan ] 07/2013-06/2018 \\
	Thesis advisor: Christine Aidala\\
	Department of Physics 
	
	\item [Undergraduate Research Assistant, University of Kentucky] 05/2012-05/2013 \\
	Thesis advisor: Renee Fatemi\\
	Department of Physics 
	
	\item [High School Research Assistant, University of Kentucky]  01/2008-05/2009 \\
	Thesis advisor: James McDonough \\
	Department of Mechanical Engineering

\end{description}





\end{flushleft}



\vspace{7pt}

\begin{flushleft}
	\LARGE\textbf{Research Funding}
	\vspace{4pt}
	\normalsize
	\begin{center}
		\begin{enumerate}
			\item Oak Ridge National Laboratory subcontract through Brookhaven National Laboratory. \textit{Software Development for sPHENIX}. 03/2020-03/2022. PI, \$466k.
		\end{enumerate}
	\end{center}
\end{flushleft}


\vspace{7pt}

\begin{flushleft}
\LARGE\textbf{Peer Reviewed Publications with Significant Contribution}
\vspace{4pt}

\small 
Full publication list at end of CV, or search `find a J. D. Osborn' on \href{inspirehep.net}{inspirehep.net} 
\begin{center}
\normalsize
\begin{enumerate}
	\item U. A. Acharya et al. ``Improving constraints on gluon spin-momentum correlations in transversely
	polarized protons via midrapidity open-heavy-flavor electrons in $p^{\uparrow}+p$ collisions at $\sqrt{s}=200$ GeV.'' arXiv:2204.12899, submitted to Phys. Rev. Lett. 
	\item J. C. Bernaeur et al. ``Scientific computing plan for the ECCE detector at the Electron Ion Collider'' arXiv:2205.08607, submitted to NIMA.
	\item X. Ai et al. ``A Common Tracking Software Project.'' Computing and Software for Big Science 6, 8 (2022).
	
	\item R. Abdul Khalek et al. ``Science Requirements and Detector Concepts for the Electron-Ion Collider: EIC Yellow Report." arXiv: 2103.05419. Accepted by Nucl. Phys. A.
	
	\item {\textbf{J.D. Osborn}} et al. ``Implementation of ACTS into sPHENIX track reconstruction.'' Computing and Software for Big Science 5, 23 (2021).
	
	\item U. A. Acharya et al. (PHENIX Collaboration), ``Probing gluon spin-momentum correlations in transversely polarized protons through midrapidity isolated direct photons in $p^\uparrow+p$ collisions at $\sqrt{s}=200$ GeV. Phys. Rev. Lett. 127, 162001 (2021).
	
	\item U. A. Acharya et al. (PHENIX Collaboration), ``Transverse single-spin asymmetries of midrapidity $\pi^0$ and $\eta$ mesons in polarized $p$$+$$p$ collisions at $\sqrt{s}=200$ GeV.'' Phys. Rev. D 103, 052009 (2021).
	
	\item C. A. Aidala et al. ``Design and beam test results for the 2D projective sPHENIX electromagnetic calorimeter prototype.'' IEEE Transactions on Nuclear Science, vol. 68, no. 2, pp. 173-181, Feb. 2021.
	
	\item J. D. Osborn for the sPHENIX Collaboration, ``Requirements, status, and plans for track reconstruction at the sPHENIX experiment.'' arXiv:2007.00771, peer reviewed Connecting the Dots 2020 Workshop proceedings.
	
	\item U. Acharya et al. (PHENIX Collaboration), ``Measurement of jet-medium interactions via direct photon-hadron correlations in Au+Au and $d+$Au collisions at $\sqrt{s_{_{NN}}}=200$ GeV.'' Phys. Rev. C 102, 054910 (2020).
	
	\item J. D. Osborn for the LHCb Collaboration, ``Jet hadronization at LHCb.'' PoS High-$p_T$ 2019 (2020) 013.
	
	\item I. Helenius, J. Lajoie, {\bf{J. D. Osborn}}, P. Paakkinen, H. Paukkunen, ``Nuclear gluons at RHIC in a multi-observable approach.'' Phys. Rev. D 100, 014004 (2019).
	
	\item A. Aaij et al. (LHCb Collaboration), ``Measurement of charged-hadron production in $Z$-tagged jets in proton-proton collisions at $\sqrt{s}=8$ TeV.'' Phys. Rev. Lett. 123, 232001 (2019).
	
	\item C. Aidala et al. (PHENIX Collaboration), ``Nonperturbative transverse momentum broadening in dihadron angular correlations in $\sqrt{s_{_{NN}}}$=~200 GeV proton-nucleus collisions.''  Phys. Rev. C 99, 044912 (2019).

	\item J. D. Osborn for the PHENIX Collaboration, ``PHENIX results on jet modification with $\pi^0$- and photon-triggered two particle correlations in $p$$+$$p$, $p(d)$+Au, and Au+Au collisions.'' Nuclear Physics A 982, 591-594 (2019).

	\item C. Aidala et al. (PHENIX Collaboration), ``Nonperturbative transverse momentum dependent effects in dihadron and direct photon-hadron angular correlations in $p$$+$$p$ collisions at $\sqrt{s}=200$ GeV.'' Phys. Rev. D 98, 072004 (2018).
		
		
	\item C. Aidala et al. (PHENIX Collaboration), ``Single-spin asymmetry of $J/\psi$ production in $p$$+$$p$, $p$+Al, and $p$$+$Au collisions with transversely polarized proton beams at $\snn$~=~200 GeV.'' Phys. Rev. D 98, 012006 (2018).
	
	\item J. D. Osborn for the PHENIX Collaboration, ``Study of cold and hot nuclear matter effects on jets with direct photon triggered correlations from PHENIX.'' Nuclear Physics A 967, 476-479 (2017).
	
	\item A. Adare et al. (PHENIX Collaboration), ``Nonperturbative-transverse-momentum effects and evolution in dihadron and direct photon-hadron angular correlations in $p$$+$$p$ collisions at $\sqrt{s}=510$ GeV.''  Phys. Rev. D95, 072002 (2017).
	
\end{enumerate}
\end{center}
\end{flushleft}




\newpage



%%%%%%%%%%%%%%%%%%%%%%%%%%%%%%%%%%%%%%%%%%%%%%%%%%%%%%%
%%%%%%%															%%%%%%%%%%
%%%%%%					AWARDS AND FELLOWSHIPS				%%%%%%%%%%
%%%%%%%															%%%%%%%%%%
%%%%%%%%%%%%%%%%%%%%%%%%%%%%%%%%%%%%%%%%%%%%%%%%%%%%%%%

\vspace{7pt}
\begin{flushleft}
	\LARGE\textbf{Awards and Fellowships} \\
	\normalsize
	
	\begin{tabular}{ll}
		Oak Ridge National Laboratory Individual Performance Award & \hspace*{3.5cm}12/2021 \\
		CSM Division Outstanding Postdoc Award & \hspace*{3.5cm}12/2021 \\
		RHIC/AGS Users' Executive Committee Merit Award & \hspace*{3.5cm}05/2019\\
		Young PHENIXian award & \hspace*{3.5cm}12/2015\\
		Michigan Graduate $1^{st}$ Year Fellowship & \hspace*{3.5cm}08/2013 \\
		Summa Cum Laude, University of Kentucky & \hspace*{3.5cm}05/2013 \\
		Sigma Pi Sigma Physics Honor Society & \hspace*{3.5cm}05/2013 \\
		University of Kentucky Outstanding Senior Award & \hspace*{3.5cm}05/2013\\
	\end{tabular}
\end{flushleft}







%%%%%%%%%%%%%%%%%%%%%%%%%%%%%%%%%%%%%%%%%%%%%%%%%%%%%%%
%%%%%%%															%%%%%%%%%%
%%%%%%					CONFERENCE ORGANIZATION				%%%%%%%%%%
%%%%%%%															%%%%%%%%%%
%%%%%%%%%%%%%%%%%%%%%%%%%%%%%%%%%%%%%%%%%%%%%%%%%%%%%%%


\vspace{7pt}
\begin{flushleft}
	\LARGE\textbf{Conference and Workshop Organization} \\
	\normalsize 
	\begin{enumerate}
	
		\item Streaming Readout 9 Workshop, virtual, December 8-10, 2021.
		\item 2021 CFNS Summer School Software Lectures, virtual, August 9-23, 2021.
		\item ECCE Second Simulation Workshop, virtual, May 21, 2021.
		\item ECCE Simulation Workshop, virtual, April 2, 2021.
		\item LHCP Conference QCD Organizing Committee, Paris, France, June 7-12, 2021.
		\item Spin and EIC workshops at the RHIC All Users Meeting, Brookhaven National Laboratory, June 2-5, 2019.
		\item Scientific Secretary, American Physical Society Division of Particles and Fields Meeting, Ann Arbor, MI, August 3-7, 2015.
	\end{enumerate}
\end{flushleft}

\vspace{7pt}



%%%%%%%%%%%%%%%%%%%%%%%%%%%%%%%%%%%%%%%%%%%%%%%%%%%%%%%%%%%%%%%%%%%%%%%%%%
%%%%%%%															%%%%%%%%%%
%%%%%%					Research Community Service				%%%%%%%%%%
%%%%%%%															%%%%%%%%%%
%%%%%%%%%%%%%%%%%%%%%%%%%%%%%%%%%%%%%%%%%%%%%%%%%%%%%%%%%%%%%%%%%%%%%%%%%%

\begin{flushleft}
	\LARGE\textbf{Collaboration and Research Community Service} \\
	\normalsize 
	\vspace*{0.3cm}
	\textbf{EIC Detector 1 Simulation and QA Convener}, 2022-present \\ Served as convenor of the simulation and QA working group for the EIC project detector transitional period post proposals.\\
	\vspace*{0.3cm}
	\textbf{JINST Reviewer} 2021-present \\ Serve as a peer reviewer for articles submitted to the Journal of Instrumentation (JINST) \\
	\vspace{0.3cm}
	\textbf{ECCE Software and Computing Convenor}, 2021-present \\ Served as convenor of the software and computing working group for the ECCE proto-collaboration.\\
	\vspace*{0.2cm}
	\textbf{sPHENIX Publication Policy Committee}, 2021-2022 \\ Served on committee developing sPHENIX publication policies and criteria. \\
	\vspace*{0.2cm}
	\textbf{sPHENIX Track Reconstruction Convenor} 2020-present \\ Served as convenor of the track reconstruction software effort for the sPHENIX collaboration. \\
	\vspace*{0.2cm}
	\textbf{sPHENIX TPOT Reviewer}, 12/2020 \\ Review committee for TPC Outer Tracker for sPHENIX. \\

	\vspace*{0.2cm}
	\textbf{DOE Office of Nuclear Physics SBIR Proposal Reviewer}, 11/2020 \\ Review SBIR proposals related to computing in Nuclear Physics. \\
	\vspace*{0.2cm}
	\textbf{Nuclear Physics Day on Capitol Hill}, 04/2018. \\ Met with staff members of Michigan Senators and Representatves to discuss funding for nuclear physics research.\\
	\vspace*{0.2cm}
	\textbf{PHENIX Executive Council}, 01/2016-01/2017.\\ Served as elected member to executive council, which advises PHENIX spokesperson on scientific priorities.
	
	
\end{flushleft}







%%%%%%%%%%%%%%%%%%%%%%%%%%%%%%%%%%%%%%%%%%%%%%%%%%%%%%%
%%%%%%%															%%%%%%%%%%
%%%%%%					SEMINARS								%%%%%%%%%%
%%%%%%%															%%%%%%%%%%
%%%%%%%%%%%%%%%%%%%%%%%%%%%%%%%%%%%%%%%%%%%%%%%%%%%%%%%

\vspace{7pt}
\begin{flushleft}
	\LARGE\textbf{Invited Conference and Workshop Presentations}
	\normalsize
	\begin{enumerate}
		\item RHIC/AGS All Users Meeting, June 7-10, 2022. ``sPHENIX Spin and Cold QCD Future Physics Program.''
		\item Connecting the Dots Workshop, May 31-June 2, 2022. ``4D track reconstruction with sPHENIX.''
		\item Streaming Readout 10 Workshop, May 17-19, 2022. ``4D track reconstruction with sPHENIX.''
		\item Probing QCD at High Energy and Density with Jets INT Workshop, August 16-20, 2021. ``Recent results and future directions for jet hadronization and substructure measurements.''
		\item CFNS EIC Summer School, August 9-23, 2021. ``ECCE software stack and detector geometry.'' 
		\item EICUG Summer Meeting, August 2-6, 2021. ``ECCE software.''
		\item RHIC Science Program Toward EIC in the Coming Years, May 24-26, 2021. ``Photon-jet and dijet probes at RHIC towards the EIC.''
		\item {\bf{Plenary}}: Computing in High Energy Physics (CHEP) Conference, May 17-21, 2021. ``Implementation of ACTS into sPHENIX Track Reconstruction.''
		\item Streaming Readout VIII Workshop, April 28-30, 2021. ``ORNL Streaming Readout Plans at the EIC.''\\
		\item Resummation, Evolution, Factorization Workshop, December 7-11, 2020. ``Experimental perspective on hadronization.''
		\item Jets for 3D Imaging Workshop, November 23-25, 2020. ``Jet substructure at the EIC.''
		\item ACTS Workshop, May 25-29, 2020. ``sPHENIX Experience with Acts.''
		\item Second EIC Yellow Report Workshop at Pavia University, May 20-22, 2020. ``Inclusive and heavy flavor jet substructure at the EIC.''
		\item Connecting The Dots Workshop, April 22-24, 2020. ``Requirements, status and plans for track reconstruction of the sPHENIX experiment.''
		\item 3$^{rd}$ JETSCAPE Workshop, March 18-20, 2020. ``Hadronization and jet substructure at RHIC and the LHC.''
		\item 13$^{th}$ International Workshop on High $p_T$ Physics in the RHIC/LHC Era, March 19-22, 2019. ``Jet hadronization at LHCb.''
		\item Workshop on Novel Probes of the Nucleon Structure in SIDIS, e+e- and pp (FF2019), March 14-16, 2019. ``Factorization breaking, color entanglement, and hadronization of jets.''
		\item Workshop on the Definition of Jets, Brookhaven National Laboratory, June 25-27, 2018. ``Probing effects from QCD color with photon-jet and dijet correlations.''
		\item CMS SMP-J Workshop Annual Workshop, January 25, 2017. ``Color entanglement and color coherence."
	\end{enumerate}
\end{flushleft}


\vspace{7pt}
\begin{flushleft}
	\LARGE\textbf{Seminars and Colloquia} \\
	\normalsize
	\begin{enumerate}
		\item University of Tennessee seminar, March 15, 2022. ``sPHENIX charged particle reconstruction.''
		\item Oak Ridge National Laboratory seminar, June 16, 2021. ``Addressing data reduction challenges with next generation scientific software.''
		\item Workflow Workshop Seminar Series, August 14, 2020. ``Data processing workflows at scattering user facilities.''
		\item University of Tennessee HEP/Nuclear/Astro Seminar, February 3, 2020. ``Jet substructure at RHIC and the LHC.''
		\item University of Kentucky HEP/Nuclear Seminar, November 7, 2019. ``Hadronization and jet substructure at RHIC and the LHC.''
		\item University of Michigan HEP/Astro/Nuclear Seminar, April 8, 2019. ``Jet substructure at RHIC and the LHC.''
		\item Oak Ridge National Laboratory Seminar, February 18, 2019. ``Peering inside protons and nuclei.''
		\item Wayne State University PAN Seminar, September 14, 2018. ``Effects from color flow in proton-proton and proton-nucleus collisions.''
		\item Paul Laurence Dunbar High School Senior Seminar, August 31, 2018. ``Career paths with a physics degree.''
		\item University of Illinois HEP/MEP Colloquium, October 23, 2017. ``Effects from color entanglement in proton-proton and proton-nucleus collisions."
		\item University of Michigan HEP/Astro/Nuclear Seminar, November 21, 2016. ``Recent experimental results on QCD factorization breaking at RHIC."
		\item Brookhaven National Laboratory Nuclear Seminar Series, October 25, 2016. ``Recent experimental results on QCD factorization breaking at RHIC."
		\item Seminar at Columbia University, October 24, 2016. ``Recent experimental results on QCD factorization breaking at RHIC."
		
	\end{enumerate}
\end{flushleft}

\vspace{7pt}










%%%%%%%%%%%%%%%%%%%%%%%%%%%%%%%%%%%%%%%%%%%%%%%%%%%%%%%
%%%%%%%															%%%%%%%%%%
%%%%%%		CONFERENCE AND WORKSHOP PRESENTATIONS			%%%%%%%%%%
%%%%%%%															%%%%%%%%%%
%%%%%%%%%%%%%%%%%%%%%%%%%%%%%%%%%%%%%%%%%%%%%%%%%%%%%%%

\begin{flushleft}
	\LARGE\textbf{Conference and Workshop Presentations} \\
	\normalsize
	\begin{enumerate}
		\item sPHENIX Summer School, May 26-27, 2022. ``Analysis and Hands on Fun4All Tutorial.''
		\item sPHENIX Collaboration Meeting, virtual, January 7, 2022. ``Tracking developments and performance.''
		\item ECCE Simulation Workshop, virtual, May 21, 2021. ``Developing an analysis package for ECCE DST analysis.''
		\item ECCE Simulation Workshop, virtual, April 2, 2021. ``Analysis in Fun4All.''
		\item sPHENIX Collaboration Meeting, Boulder, Colorado (virtual), July 6, 2020. ``ACTS Tracking.''
		\item sPHENIX Software Mega-Workfest, Brookhaven National Laboratory, January 13-17, 2020. ``Building a Fun4All Analysis Package Tutorial.''
		\item sPHENIX Collaboration Meeting, Brookhaven National Laboratory, June 5-6, 2018. ``Photons and clustering in sPHENIX.''
		\item Quark Matter 2018, Lido, Italy, May 13-19, 2018. ``PHENIX results on jet modification with $\pi^0$- and photon-triggered two particle correlations in $p$$+$$p$, $p$(d)+Au, and Au(Cu)+Au collisions.''
		
		\item PHENIX Collaboration Meeting, Brookhaven National Laboratory, December 1-3, 2017. ``High $p_T$ correlations analysis highlights and future plans.''
		
		\item RHIC Users Meeting Proton Structure Workshop, Brookhaven National Lab, June 20-23, 2017. ``Partonic structure of nucleons and nuclei at sPHENIX."
		
		\item Quark Matter 2017, Chicago, Illinois, February 5-11, 2017. ``Study of cold and hot nuclear matter effects on jets with direct photon-triggered correlations from PHENIX."
		
		
		\item 22$^{\rm nd}$ International Spin Symposium, Urbana, IL, September 25-30, 2016. ``Nonperturbative transverse momentum effects in dihadron and direct photon-hadron angular correlations."
		
		\item 4th Workshop on the QCD Structure of the Nucleon (QCD-N'16), Getxo, Spain, July 11-15, 2016. ``Nonperturbative transverse momentum effects in dihadron and direct photon-hadron angular correlations."
		
		\item American Physical Society Division of Nuclear Physics Meeting, Sante Fe, NM, October 28-31, 2015. ``Measuring intrinsic partonic transverse momentum via two-particle correlations in PHENIX."
		
		\item American Physical Society Division of Particles and Fields Meeting, Ann Arbor, MI, August 3-7, 2015. ``Parton dynamics at PHENIX." 
		
		\item Physics Graduate Student Symposium, July 8, 2015.  ``Partonic dynamics in high energy proton-proton collisions at PHENIX." 
		
		\item American Physical Society Southeastern Section Meeting, Tallahassee, FL, November 14-17, 2012.  ``Noise analysis of the Forward GEM Tracker at STAR." 
		
		\item P.L. Dunbar Math Science and Technology Research Symposium, Lexington, KY, April 26, 2009. ``Neutrinos and shock reignition in the gain region of type IIa supernovae."
	\end{enumerate}
	
\end{flushleft}

\vspace{7pt}




%%%%%%%%%%%%%%%%%%%%%%%%%%%%%%%%%%%%%%%%%%%%%%%%%%%%%%%
%%%%%%%															%%%%%%%%%%
%%%%%%				CONFERENCE PROCEEDINGS					%%%%%%%%%%
%%%%%%%															%%%%%%%%%%
%%%%%%%%%%%%%%%%%%%%%%%%%%%%%%%%%%%%%%%%%%%%%%%%%%%%%%%


\begin{flushleft}

\vspace{7pt}
\LARGE\textbf{Other Public Notes with Significant Contribution} \\
\vspace{2pt}
	\begin{center}
	\normalsize
	\begin{enumerate}
		\item sPHENIX Conceptual Design Report (CDR). The sPHENIX Collaboration. \textcolor{red}{\href{https://indico.bnl.gov/event/4640/}{https://indico.bnl.gov/event/4640}}, 2018.
		 \item Detector design study for an EIC detector around the sPHENIX solenoid. C. Aidala et al. \textcolor{red}{\href{https://indico.bnl.gov/event/5283/}{https://indico.bnl.gov/event/5283/}}, 2018.
		 \item sPHENIX medium energy barrel physics. The sPHENIX Collaboration. \textcolor{red}{\href{https://indico.bnl.gov/event/3866/}{https://indico.bnl.gov/event/3866/}}, 2017.
 		 \item T1044-2017 sPHENIX test beam EMCal analysis. J. D. Osborn and J. Huang. \textcolor{red}{\href{https://indico.bnl.gov/event/3854/}{https://indico.bnl.gov/event/3854/}}, 2017.
		\item sPHENIX modest forward upgrade LOI. The sPHENIX Collaboration. \textcolor{red}{\href{https://indico.bnl.gov/event/3867/}{https://indico.bnl.gov/event/3867}}, 2017.
		\item \textbf{Nonperturbative transverse momentum effects in dihadron and direct photon-hadron angular correlations}, J. D. Osborn for the PHENIX Collaboration. Proceedings for the SPIN 2016 Symposium, Urbana, IL, September 25-30 2016. arXiv:1701:00681.
		\item \textbf{Parton dynamics at PHENIX}, J. D. Osborn for the PHENIX Collaboration. Proceedings for the American Physical Society Division of Particles and Fields Conference, Ann Arbor, MI, August 4-8 2015. arXiv:1511.00016. 
	

\end{enumerate}

\end{center}

\end{flushleft}




\vspace{7pt}






%%%%%%%%%%%%%%%%%%%%%%%%%%%%%%%%%%%%%%%%%%%%%%%%%%%%%%%
%%%%%%%															%%%%%%%%%%
%%%%%%					TEACHING								%%%%%%%%%%
%%%%%%%															%%%%%%%%%%
%%%%%%%%%%%%%%%%%%%%%%%%%%%%%%%%%%%%%%%%%%%%%%%%%%%%%%%




\begin{flushleft}
\LARGE \textbf{Student Advising and Mentorship}\\
\normalsize
\vspace{0.3cm}
\textbf{Graduate Students Supervised}\\
Sebastian Araya - Track pattern recognition in ECCE. \\
Dillon Fitzgerald - Heavy flavor jet substructure and hadronization at the EIC. \\
Kara Mattioli - Heavy flavor jet tagging and substructure at LHCb. \\
Jordan Roth - $Z^0$-hadron and $Z^0$-jet correlations at LHCb. \\
Nicole Lewis - $\pi^0$, $\eta$, and direct photon single spin asymmetries at PHENIX. \\

\vspace{0.5cm}

\textbf{University Research Opportunity Program (UROP) at University of Michigan}, 09/2018-05/2019\\
Propose and guide a research project to an undergraduate student. Supervised Brandon Liang in a jet substructure Monte Carlo study to better understand nonperturbative contributions to jets.\\
\vspace*{0.3cm}

\textbf{Additional Undergraduate Students Supervised}\\
Nikhil Shankar, Hayden Hansen, Ezra Lesser, Nick Melekian, Ruby Araj, Emily Camras, Robert Read, Robert Cernak, Aaron White.
\end{flushleft}



\vspace{7pt}
\begin{flushleft}
\LARGE\textbf{Teaching} \\
\normalsize 
\vspace{-4pt}

\begin{description}[style=nextline]
	\item [University of Michigan] \hfill
		\vspace{-0.5cm}
		\begin{enumerate}
			\item Physics for the Life Sciences Laboratory 1: Fall 2013, Spring 2014, Fall 2014, Fall 2017
			
		\end{enumerate}
	\item [University of Kentucky] \hfill
		\vspace{-0.5cm}
		\begin{enumerate}
			\item Physics 241 - General University Physics Laboratory: Fall 2012
			\item Arts and Sciences Wired Course - Measuring Science: Fall 2011
		\end{enumerate}
\end{description}


%\begin{tabular}{lll}
%\vspace{0.5cm}
%University of Michigan & Physics for the Life Sciences Laboratory 1 & \hspace{0.35cm} \specialcell{Fall 2013\\Spring 2014\\Fall 2014} \\
%University of Kentucky & Physics 241: General University Physics Laboratory & \hspace{0.45cm}Fall 2012 \\
% & Arts and Sciences Wired Course: Measuring Science & \hspace{0.45cm}Fall 2011\\
%\end{tabular}


\end{flushleft}












%%%%%%%%%%%%%%%%%%%%%%%%%%%%%%%%%%%%%%%%%%%%%%%%%%%%%%%
%%%%%%%															%%%%%%%%%%
%%%%%%					OTHER EXPERIENCE						%%%%%%%%%%
%%%%%%%															%%%%%%%%%%
%%%%%%%%%%%%%%%%%%%%%%%%%%%%%%%%%%%%%%%%%%%%%%%%%%%%%%%



















\vspace{7pt}


































%%%%%%%%%%%%%%%%%%%%%%%%%%%%%%%%%%%%%%%%%%%%%%%%%%%%%%%
%%%%%%%															%%%%%%%%%%
%%%%%%					REFERENCES								%%%%%%%%%%
%%%%%%%															%%%%%%%%%%
%%%%%%%%%%%%%%%%%%%%%%%%%%%%%%%%%%%%%%%%%%%%%%%%%%%%%%%


%
%\begin{flushleft}
%\Large\textbf{References}\\
%\normalsize
%\vspace{0.3cm}
%
%\begin{tabular}{lll}
%\specialcell{Christine Aidala \\
%Associate Professor \\
%University of Michigan \\
%caidala@umich.edu \\
%(734) 764-1795 }
%&
%\hspace*{1.5cm}
%
%\specialcell{Jin Huang \\
%Associate Physicist \\
%Brookhaven National Laboratory \\
%jhuang@bnl.gov \\
%(631) 344-5898 
%}
%&
%\hspace*{1.5cm}
%\specialcell{Renee Fatemi \\
%Associate Professor \\
%University of Kentucky \\
%rfatemi@pa.uky.edu \\
%(859) 257-2664 \\
%}
%\end{tabular}
%
%\end{flushleft}
%








\end{document}
